\subsection*{4.3 Curl of a Vector Field}

\begin{conceptbox}
    The curl measures the tendency of a vector field to rotate:
    \begin{equation*}
        \nabla \times \mathbf{F} =
        \begin{vmatrix}
            \mathbf{i}                  & \mathbf{j}                  & \mathbf{k}                  \\
            \frac{\partial}{\partial x} & \frac{\partial}{\partial y} & \frac{\partial}{\partial z} \\
            P                           & Q                           & R
        \end{vmatrix}
    \end{equation*}
\end{conceptbox}

\begin{examplebox}
    Find the curl of $\mathbf{F}=\langle y, -x, 0 \rangle$.

    \textbf{Step 1: Identify components}
    \begin{align*}
        P=y,\quad Q=-x,\quad R=0
    \end{align*}

    \textbf{Step 2: Compute determinant}
    \begin{align*}
        \nabla \times \mathbf{F}
         & =
        \begin{vmatrix}
            \mathbf{i} & \mathbf{j} & \mathbf{k} \\
            \partial_x & \partial_y & \partial_z \\
            y          & -x         & 0
        \end{vmatrix}
    \end{align*}

    \textbf{Step 3: Expand determinant}
    \begin{align*}
        \nabla \times \mathbf{F}
         & = \mathbf{i}(0-0)
        - \mathbf{j}(0-0)
        + \mathbf{k}(-1 - 1)        \\
         & = \langle 0,0,-2 \rangle
    \end{align*}

    \textbf{Interpretation:} The field rotates about the $z$-axis.
\end{examplebox}

\begin{examplebox}

    We are given the vector field:
    \begin{equation*}
        \mathbf{F}(x,y,z) = \langle P, Q, R \rangle = \langle xz, yz, xy \rangle
    \end{equation*}

    Find curl of $\mathbf{F}=\langle xz, yz, xy \rangle$.

    \textbf{Step 1: Recall the definition of curl}

    \begin{equation*}
        \nabla \times \mathbf{F} =
        \begin{vmatrix}
            \mathbf{i}                  & \mathbf{j}                  & \mathbf{k}                  \\
            \frac{\partial}{\partial x} & \frac{\partial}{\partial y} & \frac{\partial}{\partial z} \\
            P                           & Q                           & R
        \end{vmatrix}
    \end{equation*}

    Expanding this determinant gives:
    \begin{align*}
        \nabla \times \mathbf{F}
        = \left\langle
        \frac{\partial R}{\partial y} - \frac{\partial Q}{\partial z},
        \frac{\partial P}{\partial z} - \frac{\partial R}{\partial x},
        \frac{\partial Q}{\partial x} - \frac{\partial P}{\partial y}
        \right\rangle
    \end{align*}

    \textbf{Step 2: Identify the components}

    \begin{align*}
        P = xz, \quad Q = yz, \quad R = xy
    \end{align*}

    \textbf{Step 3: Compute each component carefully}

    \textbf{First component:}
    \begin{align*}
        \frac{\partial R}{\partial y}      & = \frac{\partial}{\partial y}(xy) = x \\
        \frac{\partial Q}{\partial z}      & = \frac{\partial}{\partial z}(yz) = y \\
        \Rightarrow \text{First component} & = x - y
    \end{align*}

    \textbf{Second component:}
    \begin{align*}
        \frac{\partial P}{\partial z}       & = \frac{\partial}{\partial z}(xz) = x \\
        \frac{\partial R}{\partial x}       & = \frac{\partial}{\partial x}(xy) = y \\
        \Rightarrow \text{Second component} & = x - y
    \end{align*}

    \textbf{Third component:}
    \begin{align*}
        \frac{\partial Q}{\partial x}      & = \frac{\partial}{\partial x}(yz) = 0 \\
        \frac{\partial P}{\partial y}      & = \frac{\partial}{\partial y}(xz) = 0 \\
        \Rightarrow \text{Third component} & = 0 - 0 = 0
    \end{align*}

    \textbf{Step 4: Combine results}

    \begin{align*}
        \nabla \times \mathbf{F} = \langle x - y,\ x - y,\ 0 \rangle
    \end{align*}

    \textbf{Final Answer:}
    \begin{equation*}
        \nabla \times \mathbf{F} = \langle x - y,\ x - y,\ 0 \rangle
    \end{equation*}

\end{examplebox}
\begin{conceptbox}
    Divergence and curl provide physical meaning:
    \begin{itemize}
        \item Divergence: source or sink strength
        \item Curl: rotational tendency
    \end{itemize}
\end{conceptbox}

\begin{examplebox}
    Fluid flow interpretation:
    \begin{itemize}
        \item If $\nabla \cdot \mathbf{v} > 0$: fluid expands (source)
        \item If $\nabla \cdot \mathbf{v} < 0$: fluid compresses (sink)
    \end{itemize}
    Rotation example:

    A whirlpool has nonzero curl, indicating rotation in the flow field.
\end{examplebox}
\begin{conceptbox}
    A vector field is:
    \begin{itemize}
        \item Conservative if $\nabla \times \mathbf{F} = 0$
        \item Solenoidal if $\nabla \cdot \mathbf{F} = 0$
    \end{itemize}
\end{conceptbox}

\begin{examplebox}
    Check if $\mathbf{F}=\langle 2x,2y,2z\rangle$ is conservative.

    \begin{align*}
        \nabla \times \mathbf{F} = \mathbf{0}
    \end{align*}

    Thus, it is conservative.
\end{examplebox}

\begin{examplebox}
    Check if $\mathbf{F}=\langle -y,x,0\rangle$ is solenoidal.

    \begin{align*}
        \nabla \cdot \mathbf{F} = 0
    \end{align*}

    Thus, it has no sources or sinks.
\end{examplebox}